\documentclass{article}
\usepackage{graphicx} % Required for inserting images
\usepackage{amsmath}

\title{tecniche fusion}
\date{June 2026}

\begin{document}

\maketitle

\section{Per-CWE Analysis}

\section{Behavior-Knowledge Space (BKS)}
Let $L$ be the number of analyzers and let $e_i(x)\in\{0,1\}$ denote the decision produced by analyzer $i$ for a code artifact $x$,
where $1$ indicates the detection of a vulnerability and $0$ otherwise.
Each artifact is therefore represented by a behavior vector

\begin{equation}
E(x) = (e_1(x), e_2(x), \ldots, e_L(x)),
\end{equation}

which encodes the joint behavior of all analyzers.
During training, the Behavior-Knowledge Space (BKS) stores, for each observed behavior vector $E$,
the number of occurrences associated with each ground-truth class $c$.
Let $N(E,c)$ denote the number of training instances exhibiting behavior pattern $E$ and belonging to class $c$.
The posterior probability of class $c$ given the observed behavior is estimated empirically as

\begin{equation}
P(c \mid E) =
\frac{N(E,c)}
{\sum_{c'} N(E,c')}.
\end{equation}

During inference, a new artifact is assigned to the class maximizing the estimated posterior probability:

\begin{equation}
\hat{c} =
\arg\max_c P(c \mid E).
\end{equation}

(Per CWE analysis): As an illustrative example, consider three static analyzers, namely CodeQL, Semgrep, and Joern.
Suppose a code fragment produces the behavior vector $E=(1,0,0)$, indicating that CodeQL reports a vulnerability while Semgrep and Joern do not.
If the training set contains $50$ occurrences of this pattern, of which $42$ correspond to true vulnerabilities
and $8$ to non-vulnerable artifacts, the estimated posterior probability is

\begin{equation}
P(\text{vulnerable} \mid E=(1,0,0))
=
\frac{42}{42+8}
=
0.84.
\end{equation}

Consequently, the artifact is classified as vulnerable with an estimated confidence of $84\%$.
Unlike majority-voting schemes which only consider the number of analyzers in agreement (and would have set this combination as safe),
BKS explicitly models the empirical relationship between specific combinations of findings and the corresponding ground truth.
Furthermore, unlike Na\"ive Bayes-based fusion approaches, BKS does not require conditional independence among analyzer outputs,
allowing it to capture the strong dependencies that typically arise between static analysis tools operating on the same source code.
This property is particularly valuable in vulnerability detection, where recurring patterns of agreement and disagreement
among analyzers often provide stronger predictive power than individual tool decisions considered in isolation.


\section{Naive Bayes}

Let $L$ be the number of analyzers and let $e_i(x)\in\{0,1\}$ denote the binary decision produced by analyzer $i$ for a code artifact $x$,
where $1$ indicates the detection of a vulnerability and $0$ otherwise.
Given the behavior vector

\begin{equation}
E(x) = (e_1(x), e_2(x), \ldots, e_L(x)),
\end{equation}

Naive Bayes estimates the posterior probability of each class $c$ using Bayes' theorem.
Assuming conditional independence among analyzer outputs given the class label,
the posterior probability can be expressed as

\begin{equation}
P(c \mid E)
=
\frac{P(c)\prod_{i=1}^{L}P(e_i \mid c)}
{\sum_{c'} P(c')\prod_{i=1}^{L}P(e_i \mid c')}.
\end{equation}

The prior probability $P(c)$ and the conditional probabilities $P(e_i \mid c)$ are estimated from the training data.
During inference, a new artifact is assigned to the class maximizing the posterior probability:

\begin{equation}
\hat{c}
=
\arg\max_c P(c \mid E).
\end{equation}

(Per CWE analysis): As an illustrative example, consider three static analyzers, namely CodeQL, Semgrep, and Joern.
Suppose a code fragment produces the behavior vector $E=(1,0,0)$, indicating that CodeQL reports a vulnerability while Semgrep and Joern do not.
Assume that the training data provide the following estimates:

\begin{equation}
P(\text{vulnerable}) = 0.40,
\qquad
P(\text{safe}) = 0.60,
\end{equation}

\begin{equation}
P(e_1=1 \mid \text{vulnerable}) = 0.80,
\qquad
P(e_1=1 \mid \text{safe}) = 0.15,
\end{equation}

\begin{equation}
P(e_2=0 \mid \text{vulnerable}) = 0.30,
\qquad
P(e_2=0 \mid \text{safe}) = 0.85,
\end{equation}

\begin{equation}
P(e_3=0 \mid \text{vulnerable}) = 0.40,
\qquad
P(e_3=0 \mid \text{safe}) = 0.90.
\end{equation}

The corresponding likelihoods are therefore

\begin{equation}
P(E \mid \text{vulnerable})
=
0.80 \times 0.30 \times 0.40
=
0.096,
\end{equation}

\begin{equation}
P(E \mid \text{safe})
=
0.15 \times 0.85 \times 0.90
=
0.11475.
\end{equation}

Applying Bayes' theorem yields

\begin{equation}
P(\text{vulnerable} \mid E)
=
\frac{0.40 \times 0.096}
{0.40 \times 0.096 + 0.60 \times 0.11475}
\approx
0.36.
\end{equation}

Consequently, the artifact is classified as safe, since the posterior probability of the vulnerable class is lower than that of the safe class.
Unlike majority-voting schemes, Naive Bayes incorporates both the prior prevalence of vulnerabilities and the historical reliability of individual analyzers.
However, \textbf{its effectiveness relies on the conditional independence assumption among analyzer outputs}.
In vulnerability detection scenarios, static analyzers often exhibit correlated behaviors because they inspect similar code patterns and security properties.
As a result, \textbf{the independence assumption may not always hold}, potentially limiting the ability of Naive Bayes to fully exploit recurring combinations of agreements and disagreements among analyzers.

\end{document}
